\section{Common domination and the information in exact positions}
\label{sec:v27-domination}
\label{sec:v26-domination}

Three properties must be separated: domination of a statistical family by
one sigma-finite measure, absolute continuity with respect to a designated
reference member, and mutual absolute continuity of the members.  Moving
support need not prevent the first property.  The distinction also separates
the scalar-coordinate local experiments from the richer position protocol.

\begin{proposition}[Common reference measures]
\label{prop:v26-common-domination}
The family \eqref{eq:v22-thinned-family} is dominated, for every \(n\), by
\(\mu=\operatorname{Leb}|_U+\delta_\dagger\).  Its censored family is
dominated by \(\operatorname{Leb}|_U+\delta_\dagger+\delta_\partial\).
The finite products have the corresponding product dominating measures.
The original members need not be absolutely continuous with respect to
\(P_{n,0}\); the common-collar construction establishes this stronger
property for the censored members and their censored reference.

On the fixed finite strip \(B=D\times[-R,R]\), let
\[
 d\Lambda_z=\rho(u,v)\mathbf1_{\{y>U(u,v)z\}}\,du\,dv\,dy,
 \qquad d\Lambda_* =\rho(u,v)\,du\,dv\,dy.
\]
Then the Poisson law \(\mathsf P_* =\operatorname{PPP}(\Lambda_*)\) is a
common dominating probability measure for the entire finite-strip family,
and
\begin{equation}\label{eq:v26-poisson-envelope}
 \frac{d\mathsf P_z}{d\mathsf P_*}(N)
 =\exp\{\Lambda_*(B)-\Lambda_z(B)\}
           \prod_{q\in N}\mathbf1_{\{y(q)>U(q)z\}}.
\end{equation}
The empty product is one.  Multiplication by a fixed batch proportion gives
the same assertion for each batch of \eqref{eq:v23-et-poisson-intensity}.
\end{proposition}

\begin{proof}
The first statements follow by writing the displayed Lebesgue densities
and atomic masses.  They do not involve division by the reference density.
For example, on \([0,2]\),
\[
 f_\theta(x)=\frac{2(1+\theta-x)_+}{(1+\theta)^2},\qquad 0\le\theta<1,
\]
is a family of probability densities, but
\(P_\theta((1,1+\theta))=\theta^2/(1+\theta)^2>0\) for \(\theta>0\),
whereas \(P_0((1,1+\theta))=0\).

For the Poisson assertion write \(A_z=\{y>Uz\}\cap B\).  Under
\(\mathsf P_*\), restrictions to \(A_z\) and its complement are
independent Poisson processes.  The probability of no point in the
complement is \(\exp\{-\Lambda_*(B\setminus A_z)\}\).  Multiplying its
indicator by the reciprocal of that probability has expectation one and
leaves the independent process on \(A_z\) unchanged.  The resulting law
has intensity \(\Lambda_*|_{A_z}=\Lambda_z\), proving
\eqref{eq:v26-poisson-envelope}.  The argument includes the zero-intensity
and empty-configuration cases.
\end{proof}

A zero-shift Poisson member need not dominate the others.  If \(Uz<0\) on
a set of positive weighted endpoint measure, there is positive intensity
in \(\{Uz<y<0\}\), where the reference has none.  The event of at least
one point there has positive alternative probability and zero reference
probability.  This is a failure of reference absolute continuity, not of
common domination.  It also excludes a Gaussian shift with mutually
absolutely continuous members as the local limit in such a direction.
The common-collar Gaussian limit and the moving-ceiling Poisson limit remain
different experiments for the reasons supplied by their boundary rates and
support events, not by an incorrect classification of their dominating
measures.

For a fixed compact parameter set, choose \(R\) beyond all displacements.
The positive tail outside the strip is parameter independent and the negative
tail is empty.  Adjoining the common positive-tail process preserves the
common domination just proved, also for the projective experiment of
Remark~\ref{rem:v24-projective-poisson}.  This makes no uniform assertion
about an unbounded parameter set without a common displacement bound.

\subsection{Matched records and the observed contact type}

We use the convention
\(\|P-Q\|_{\rm TV}=\sup_A|P(A)-Q(A)|\), so a deficiency belongs to
\([0,1]\).  For experiments \(\mathsf E=(P_\lambda)_{\lambda\in\Theta}\)
and \(\mathsf F=(Q_\lambda)_{\lambda\in\Theta}\), write
\[
 \delta(\mathsf E,\mathsf F)=\inf_K\sup_{\lambda\in\Theta}
                   \|KP_\lambda-Q_\lambda\|_{\rm TV},\qquad
 \Delta(\mathsf E,\mathsf F)=
       \max\{\delta(\mathsf E,\mathsf F),\delta(\mathsf F,\mathsf E)\},
\]
where the kernels are independent of the unknown parameter.

Fix an even-flight design starting at contact type \(b\).  Both retained
endpoints then have type \(b\).  In this subsection a successful position
record is exactly \((X_{b,\lambda}(u),X_{b,\lambda}(v))\), and its scalar
coarsening is exactly \((u,v)\).  Residual time is integrated out in both
records; intermediate collisions and waiting counts occur in neither.
The fixed transverse projection is one parameter-independent kernel.  All
claims concern these matched records, conditional on success.  They do not
remove preparation costs from the unconditional acquisition theorems.

\begin{lemma}[Dominated input and disjoint output supports]
\label{lem:v27-dominated-disjoint}
Suppose \(\Theta\) is uncountable, \(P_\lambda\ll\nu\) for one probability
measure \(\nu\), and there are pairwise disjoint measurable sets \(A_\lambda\)
with \(Q_\lambda(A_\lambda)=1\).  Then
\(\delta((P_\lambda),(Q_\lambda))=1\).
\end{lemma}

\begin{proof}
For any kernel \(K\), put \(Q=K\nu\).  If \(Q(A)=0\), nonnegativity
gives \(K(x,A)=0\) for \(\nu\)-almost every \(x\), hence
\(KP_\lambda(A)=0\) for every \(\lambda\).  Thus \(KP_\lambda\ll Q\).
A probability measure gives positive mass to at most countably many
pairwise disjoint measurable sets.  Choose \(\lambda\) outside that
countable set for \(Q\).  Then \(KP_\lambda(A_\lambda)=0\), whereas
\(Q_\lambda(A_\lambda)=1\).  The worst-case error of \(K\) is therefore
one.  Infimizing over \(K\) proves the assertion.  No integration over the
uncountable parameter set is used.
\end{proof}

\begin{proposition}[Observed-type dichotomy for anchored homotheties]
\label{prop:v27-observed-type-dichotomy}
\label{prop:v26-position-deficiency}
Let \(\Theta\) be a nondegenerate compact interval of positive homothety
parameters.  In a persistent registered analytic facing channel let the
obstacle of type \(b_*\) vary as
\(K_\lambda=p+\lambda(K-p)\), with \(p\in\partial K\) its fixed contact
point and fixed tangent.  Keep the facing obstacle fixed, with positive
anchored gap.  Choose one admissible even-flight design and common small
contact windows.  The successful endpoint law has positive smooth flux
and a density with respect to the two boundary arclength coordinates.
Let \(\mathsf E_{\rm sc,b}^{(k)}\) and
\(\mathsf E_{\rm pos,b}^{(k)}\) denote \(k\ge1\) independent matched
successful records of the preceding definition.

If the design starts at \(b=b_*\), the distinct position laws are mutually
singular and admit no common sigma-finite dominating measure.  For every
finite \(k\ge1\),
\begin{equation}\label{eq:v26-position-deficiency}
 \delta(\mathsf E_{\rm sc,b_*}^{(k)},\mathsf E_{\rm pos,b_*}^{(k)})=1,
 \qquad
 \delta(\mathsf E_{\rm pos,b_*}^{(k)},\mathsf E_{\rm sc,b_*}^{(k)})=0.
\end{equation}
If instead the design starts at the fixed obstacle, \(b=1-b_*\), then
\begin{equation}\label{eq:v27-fixed-type-equivalence}
        \Delta(\mathsf E_{\rm sc,b}^{(k)},\mathsf E_{\rm pos,b}^{(k)})=0.
\end{equation}
In this case the position family has a common dominating measure and a
common positive-density region, and distinct members are not mutually
singular.  In either case the scalar laws are continuous in total variation
as \(\lambda\) varies.
\end{proposition}

\begin{proof}
First suppose the observed type is \(b_*\).  Along each ray from \(p\)
entering the strictly convex body \(K\), there is exactly one other boundary
point.  A homothety multiplies its distance from \(p\) by \(\lambda\).
Thus two distinct homothetic boundaries meet only at \(p\).  Removing the
anchor gives pairwise disjoint Borel curves
\(\Gamma_\lambda=\partial K_\lambda\setminus\{p\}\).
The first endpoint has an arclength density and hits \(p\) with probability
zero.  The event that the first recorded position lies in
\(\Gamma_\lambda\) has probability one under that member, also for any
finite product of successful records.  These events are pairwise disjoint.
They prove mutual singularity.  A sigma-finite measure has at most
countably many disjoint measurable sets of positive measure: intersect
with a countable finite-measure cover, and then group by a positive lower
bound on the intersection measure.  A common dominating measure is
therefore impossible.

In the fixed transverse coordinates all scalar laws have densities on one
bounded box \(U\subset\mathbb R^2\).  They are dominated by normalized
Lebesgue measure on \(U\); their \(k\)-fold products are dominated by its
\(k\)-fold product.  Apply Lemma~\ref{lem:v27-dominated-disjoint} to the
events just constructed.  This gives the first equality in
\eqref{eq:v26-position-deficiency}.  The fixed transverse projection gives
the reverse deficiency zero.  In particular the Le Cam distance is one.

Now suppose \(b=1-b_*\).  The observed contact arc is a single fixed graph
\(X_b\), independent of \(\lambda\).  The maps
\[
 F_b(u,v)=(X_b(u),X_b(v)),\qquad
 T_b(x,x')=(\text{transverse}(x),\text{transverse}(x'))
\]
are measurable and mutually inverse on the supported contact boxes.
Moreover \(P_{\lambda,b}^{\rm pos}=(F_b)_*P_{\lambda,b}^{\rm sc}\).
Extend the maps arbitrarily off those boxes.  They are
parameter-independent kernels and give
\eqref{eq:v27-fixed-type-equivalence}, also on products.  The inverse may
use the fixed obstacle: it is a known constant of this specified family,
not an unknown parameter-dependent boundary.

For either observed type the scalar density is
\begin{equation}\label{eq:v27-scalar-finite-density}
 f_{\lambda,b}(u,v)=Z_{\lambda,b}^{-1}A_{\lambda,b}(u,v)
                  (d-E_{J,\lambda,b}(u,v))_+,
 \qquad A_{\lambda,b}>0.
\end{equation}
The action and flux vary smoothly on the common box, the normalizers are
positive, and compactness bounds them away from zero.  Uniform continuity
and \(|x_+-y_+|\le|x-y|\) prove total-variation continuity.  Since
\(E_{J,\lambda,b}(0,0)=0\), a sufficiently small common square has
\(E_{J,\lambda,b}<d/4\) for every \(\lambda\).  On this square all
scalar densities are positive.  For the fixed observed obstacle,
\((F_b)_*\operatorname{Leb}|_U\) is a common dominating measure, with a
common positive-density region.  No two position members can be mutually
singular.  The same continuity argument applies to the limiting factorized
endpoint law.
\end{proof}

\begin{remark}[Matched residual-time retention]
\label{rem:v27-matched-time}
The same dichotomy holds if residual time is retained in both records.
Use \((u,v,r)\) and \((X_{b,\lambda}(u),X_{b,\lambda}(v),r)\),
respectively.  Scalar densities are dominated on a fixed bounded
three-dimensional box.  For a fixed observed obstacle use
\(F_b\times\mathrm{id}\); for a varying observed obstacle use the same
first-position disjoint events.  Total-variation continuity follows by
integrating the difference of ceiling indicators, whose integral is
bounded by the ceiling displacement.  The common square above, together
with \(d/4<r<d/2\), is a positive-density core.  Comparing a time-retaining
record with a time-deleted record is a different question and is not an
instance of \eqref{eq:v27-fixed-type-equivalence}.
\end{remark}

\subsection{A fixed-obstacle design and nonsymmetric examples}

The dependence on the observed type already occurs for two circular
contacts.  Set
\[
 K_0=\{(x,y):(x+1)^2+y^2\le1\},\qquad
 K_{1,\lambda}=\{(x,y):(x-1-\lambda)^2+y^2\le\lambda^2\},
 \quad \lambda\in[1-\epsilon,1+\epsilon].
\]
Their closest points are \((0,0)\) and \((1,0)\), and their gap is one.
The second disk is obtained by homothety about \((1,0)\).  For small
\(\epsilon>0\), take \(J=2\), physical time \(2+d\), and the word
\(K_0\longrightarrow K_{1,\lambda}\longrightarrow K_0\).  Both observed
positions belong to the fixed arc
\[
             X_0(u)=(-1+\sqrt{1-u^2},u).
\]
The finite stationarity equations have a nondegenerate solution at the
normal orbit and hence a common smooth solution on a small box.  Its
excess vanishes at the origin.  Formula~\eqref{eq:v27-scalar-finite-density}
therefore has the common positive core used in the proof, and \(F_0\)
gives exact endpoint-only equivalence.  Unobserved intermediate visits to
the varying disk do not turn these endpoint-position laws into singular
laws.  This construction explains why the type qualification in
\eqref{eq:v26-position-deficiency} cannot be omitted.

The varying-type conclusion does not require circular symmetry.  A
strictly convex analytic oval with support function
\[
 h(\theta)=1+\tfrac1{50}\cos(2\theta)+\tfrac1{100}\sin(3\theta)
\]
has \(h+h''\ge43/50>0\).  Its nonzero harmonic orders two and three
exclude a nonidentity rotational symmetry.  Fix a boundary anchor and use
nearby homothetic copies, with a facing obstacle kept at positive gap.
A sufficiently small parameter interval preserves the local channel.
Starting on the oval gives \eqref{eq:v26-position-deficiency}; starting on
the fixed facing obstacle gives \eqref{eq:v27-fixed-type-equivalence}.
These are local examples, without an assertion of the global cycle and
spanning-tree hypotheses for an unspecified periodic table.

Exact injectivity of the transverse \emph{law} map in
Theorem~\ref{thm:v26-single-offset-global} is compatible with maximal
finite-sample deficiency on the varying-type family.  The supremum over
an uncountable parameter interval and the finite-sample record are essential.
Distribution-level determination, recovery of a particular homothety
parameter from an exact position, and reconstruction of an arbitrary
analytic obstacle are distinct assertions.  The global law inverse remains
a statement using both contact types and the separately retained gap.
